% Môn: Vật lý
% Lớp: 10
% Chương: Chương I. Mở đầu
% Bài: L10C1 Bài 2. Các quy tắc an toàn trong phòng thực hành Vật lí · Quy trình sơ cứu khẩn cấp khi xảy ra tai nạn điện giật.
% Số câu: 185
% App đọc file này trực tiếp — sửa rồi Commit trên GitHub, bấm Đồng bộ GitHub .tex

% L10C1 Bài 2. Các quy tắc an toàn trong phòng thực hành Vật lí



% ===== Câu 22 =====
% ID: L10C1B2-100-TN
% Mức: TH
\dangbt{Quy trình sơ cứu khẩn cấp khi xảy ra tai nạn điện giật.}
\begin{ex}
% ID: L10C1B2-100-TN
% Mức: TH
Tại sao người cứu hộ phải đứng trên một tấm gỗ khô, thảm cao su hoặc đi dép nhựa khô khi thực hiện các thao tác tách nạn nhân ra khỏi nguồn điện hở khi chưa ngắt được nguồn điện tổng?
\choice
{Để tăng ma sát giúp người cứu đứng vững vàng hơn khi kéo nạn nhân}
{\True Để cách điện với đất, ngăn không cho dòng điện truyền qua cơ thể người cứu xuống đất gây điện giật}
{Để giữ ấm cho lòng bàn chân không bị nhiễm lạnh từ sàn phòng thực hành}
{Để làm giảm bớt nhiệt độ tỏa ra từ hệ thống thiết bị đang chập cháy xung quanh}
\loigiai{
Việc đứng trên vật cách điện khô nhằm tạo ra một điện trở cách điện vô cùng lớn giữa cơ thể người cứu và đất, đảm bảo dòng điện không thể chạy qua người cứu xuống đất nếu vô tình chạm phải nguồn điện hở trong quá trình cứu hộ.
}
\end{ex}

% ===== Câu 23 =====
% ID: L10C1B2-101-TN
% Mức: TH
\begin{ex}
% ID: L10C1B2-101-TN
% Mức: TH
Tại sao khi xảy ra đám cháy do cồn đổ ra mặt bàn thí nghiệm, người làm thí nghiệm tuyệt đối không được dùng nước để dập lửa?
\choice
{Vì nước phản ứng hóa học mạnh mẽ với cồn sinh ra khí độc gây ngạt thở}
{\True Vì cồn nhẹ hơn nước và tan vô hạn trong nước, dập nước sẽ khiến cồn nổi lên chảy loang rộng đám cháy hơn}
{Vì nước thông thường chứa nhiều tạp chất dẫn điện làm tăng nguy cơ giật điện từ ngọn lửa cồn}
{Vì nước làm tăng nhiệt độ của đám cháy cồn khiến lửa bùng phát mạnh hơn}
\loigiai{
Cồn có khối lượng riêng nhẹ hơn nước và tan vô hạn trong nước. Khi tạt nước vào đám cháy cồn, nước không dập tắt được lửa mà hòa tan cồn rồi kéo cồn đang cháy loang rộng ra diện tích lớn hơn xung quanh bàn, làm đám cháy bùng phát dữ dội và nguy hiểm hơn.
}
\end{ex}

% ===== Câu 24 =====
% ID: L10C1B2-103-TN
% Mức: TH
\begin{ex}
% ID: L10C1B2-103-TN
% Mức: TH
Trong một thí nghiệm liên quan đến việc sử dụng nguồn nhiệt (ví dụ: đèn cồn, bếp điện), học sinh cần tuân thủ những quy tắc an toàn nào để phòng tránh hỏa hoạn và bỏng?
\choice
{Để các vật liệu dễ cháy như giấy, dung dịch dễ bay hơi gần nguồn nhiệt.}
{\True Chỉ sử dụng nguồn nhiệt khi có sự giám sát trực tiếp của giáo viên hoặc người có trách nhiệm.}
{Thổi tắt đèn cồn bằng miệng để tiết kiệm thời gian.}
{Để dung dịch hóa chất tràn ra ngoài khu vực thí nghiệm nếu không may xảy ra.}
\loigiai{
Sử dụng nguồn nhiệt trong phòng thực hành luôn tiềm ẩn nguy cơ cháy nổ và bỏng. Việc chỉ sử dụng nguồn nhiệt khi có sự giám sát của người có chuyên môn là quy tắc an toàn cơ bản nhất để xử lý kịp thời các tình huống bất ngờ. Để vật liệu dễ cháy gần nguồn nhiệt, thổi tắt đèn cồn bằng miệng, hoặc để hóa chất tràn ra ngoài đều là những hành vi nguy hiểm, vi phạm nghiêm trọng quy tắc an toàn.
}
\end{ex}

% ===== Câu 25 =====
% ID: L10C1B2-104-TN
% Mức: VD
\begin{ex}
% ID: L10C1B2-104-TN
% Mức: VD
Trong quy trình cứu hộ nạn nhân bị tai nạn điện giật, sau khi ngắt điện hoặc tách được nạn nhân ra khỏi nguồn điện một cách an toàn, nếu thấy nạn nhân đã ngừng thở thì cần thực hiện ngay lập tức thao tác sơ cứu nào tại chỗ?
\choice
{Dội nước lạnh lên mặt nạn nhân để nạn nhân tỉnh lại}
{Đưa nạn nhân vào phòng điều hòa nhiệt độ thấp để hạ thân nhiệt}
{\True Tiến hành hà hơi thổi ngạt kết hợp ép tim ngoài lồng ngực liên tục, đồng thời gọi cấp cứu y tế}
{Xoa dầu nóng vào lòng bàn tay, bàn chân và để nạn nhân nằm im tĩnh dưỡng}
\loigiai{
Khi nạn nhân ngừng thở do điện giật, việc khôi phục tuần hoàn và hô hấp là ưu tiên sống còn. Cần tiến hành hà hơi thổi ngạt phối hợp ép tim ngoài lồng ngực liên tục cho đến khi nạn nhân tự thở lại hoặc có nhân viên y tế tiếp quản.
}
\end{ex}

% ===== Câu 26 =====
% ID: L10C1B2-107-TN
% Mức: VD
\begin{ex}
% ID: L10C1B2-107-TN
% Mức: VD
Trong quá trình thực hiện thí nghiệm đốt nóng bằng đèn cồn hoặc đèn khí gas, một học sinh ngửi thấy mùi khí gas nồng nặc trong phòng thực hành. Hành động đúng đắn và cấp thiết nhất mà học sinh đó cần thực hiện ngay lập tức để đảm bảo an toàn là gì?
\choice
{Nhanh chóng tìm nguồn rò rỉ khí và cố gắng tự khắc phục.}
{\True Ngay lập tức khóa van cấp khí chính của đèn hoặc hệ thống, sau đó thông báo cho giáo viên hoặc người phụ trách phòng thí nghiệm.}
{Mở tất cả các cửa sổ và cửa ra vào để thông gió, đồng thời bật quạt thông gió.}
{Tiếp tục quan sát thí nghiệm và chờ đợi mùi khí tự bay đi.}
\loigiai{
Khi phát hiện mùi khí gas nồng nặc, điều đầu tiên và quan trọng nhất là phải loại bỏ nguồn rò rỉ khí để ngăn chặn nguy cơ cháy nổ hoặc ngộ độc. Do đó, hành động cấp thiết nhất là ngay lập tức khóa van cấp khí chính. Sau đó, việc thông báo cho giáo viên hoặc người phụ trách là bắt buộc để họ có biện pháp xử lý chuyên nghiệp và đảm bảo an toàn cho toàn bộ phòng thí nghiệm. Việc tự ý tìm và khắc phục nguồn rò rỉ có thể gây nguy hiểm hơn, và việc chỉ thông gió mà không khóa nguồn khí là không đủ để loại bỏ nguy cơ. Bật quạt thông gió cũng có thể tạo ra tia lửa điện, gây cháy nổ nếu nồng độ khí gas đủ cao.
}
\end{ex}

% ===== Câu 174 =====
% ID: L10C1B2-86-TN
% Mức: NB
\begin{ex}
% ID: L10C1B2-86-TN
% Mức: NB
Vì sao không được dùng nước để dập cháy khi thiết bị điện vẫn đang có điện?
\choice
{Vì nước làm ngọn lửa cháy sáng hơn trong mọi trường hợp}
{Vì nước làm thiết bị nhẹ hơn}
{\True Vì nước có thể dẫn điện và gây điện giật}
{Vì nước làm giảm áp suất không khí}
\loigiai{
Nước có thể dẫn điện, nên dội nước vào thiết bị đang có điện có thể gây điện giật cho người chữa cháy.
}
\end{ex}

% ===== Câu 175 =====
% ID: L10C1B2-87-TN
% Mức: NB
\begin{ex}
% ID: L10C1B2-87-TN
% Mức: NB
Khi phát hiện mùi khét từ ổ cắm điện trong phòng thực hành, việc cần làm là gì?
\choice
{Tiếp tục sử dụng để hoàn thành thí nghiệm}
{Dùng tay kiểm tra ổ cắm}
{\True Ngắt điện nếu an toàn và báo ngay cho giáo viên}
{Che ổ cắm bằng giấy}
\loigiai{
Mùi khét có thể là dấu hiệu chập điện hoặc quá tải. Cần ngắt điện nếu an toàn và báo giáo viên xử lí.
}
\end{ex}

% ===== Câu 176 =====
% ID: L10C1B2-88-TN
% Mức: NB
\begin{ex}
% ID: L10C1B2-88-TN
% Mức: NB
Vật liệu nào sau đây tuyệt đối không được sử dụng để tiếp xúc hoặc đẩy dây dẫn đang hở điện ra khỏi cơ thể người bị nạn điện giật?
\choice
{Thanh tre khô}
{Cán chổi nhựa khô}
{\True Thanh sắt ẩm ướt hoặc thanh kim loại}
{Găng tay cao su dày chịu điện chuyên dụng}
\loigiai{
Thanh sắt hay kim loại là chất dẫn điện rất tốt, đồng thời nước ẩm ướt cũng dẫn điện. Sử dụng các vật liệu này sẽ truyền dòng điện trực tiếp sang cơ thể người cứu hộ dẫn đến tai nạn điện giật kép.
}
\end{ex}

% ===== Câu 177 =====
% ID: L10C1B2-90-TN
% Mức: NB
\begin{ex}
% ID: L10C1B2-90-TN
% Mức: NB
Khi xảy ra cháy nhỏ do thiết bị điện, cách xử lí nào sau đây là phù hợp?
\choice
{Dội nước trực tiếp vào thiết bị đang có điện}
{\True Ngắt nguồn điện và dùng bình chữa cháy phù hợp}
{Dùng tay kéo thiết bị ra khỏi ổ cắm khi tay ướt}
{Đổ thêm cồn để ngọn lửa tắt nhanh}
\loigiai{
Với cháy liên quan đến điện, cần ngắt nguồn điện trước rồi dùng bình chữa cháy phù hợp, không dội nước khi thiết bị còn điện.
}
\end{ex}

% ===== Câu 178 =====
% ID: L10C1B2-91-TN
% Mức: NB
\begin{ex}
% ID: L10C1B2-91-TN
% Mức: NB
Khi đang tiến hành thực hành thí nghiệm Vật lí, nếu phát hiện xảy ra sự cố cháy do chập điện trên hệ thống dây dẫn, hành động khẩn cấp đầu tiên cần thực hiện là gì?
\choice
{Dùng vòi nước áp lực cao phun trực tiếp vào vị trí đang phát hỏa}
{\True Lập tức ngắt cầu dao điện tổng của phòng thực hành để cô lập nguồn điện}
{Chạy ra ngoài phòng tìm kiếm các phương tiện cứu hỏa chuyên dụng}
{Di tản toàn bộ học sinh và dụng cụ thí nghiệm ra khỏi phòng trước khi dập lửa}
\loigiai{
Ngắt cầu dao điện tổng là bước tiên quyết và bắt buộc để ngắt dòng điện chạy qua vị trí chập cháy, ngăn chặn tác nhân sinh nhiệt tiếp diễn và đảm bảo an toàn cho người dập lửa tránh bị điện giật.
}
\end{ex}

% ===== Câu 179 =====
% ID: L10C1B2-92-TN
% Mức: NB
\begin{ex}
% ID: L10C1B2-92-TN
% Mức: NB
Nguyên nhân nào sau đây có thể gây cháy trong phòng thực hành vật lí?
\choice
{Tắt nguồn điện sau khi thực hành}
{Sắp xếp dụng cụ gọn gàng}
{\True Để vật dễ cháy gần nguồn nhiệt}
{Đọc kĩ hướng dẫn trước khi làm thí nghiệm}
\loigiai{
Vật dễ cháy đặt gần nguồn nhiệt có thể bắt lửa, đây là một nguyên nhân gây cháy thường gặp trong phòng thực hành.
}
\end{ex}

% ===== Câu 180 =====
% ID: L10C1B2-94-TN
% Mức: NB
\begin{ex}
% ID: L10C1B2-94-TN
% Mức: NB
Khi phát hiện dây dẫn điện bị hở trong lúc thực hành, học sinh cần làm gì trước tiên?
\choice
{Tiếp tục thí nghiệm nếu mạch vẫn hoạt động}
{Dùng tay xoắn lại phần dây bị hở}
{\True Báo ngay cho giáo viên và ngừng sử dụng thiết bị}
{Đổ nước vào chỗ dây bị hở}
\loigiai{
Dây dẫn bị hở có thể gây điện giật hoặc chập điện. Cần ngừng sử dụng và báo giáo viên để xử lí an toàn.
}
\end{ex}

% ===== Câu 181 =====
% ID: L10C1B2-95-TN
% Mức: NB
\begin{ex}
% ID: L10C1B2-95-TN
% Mức: NB
Biện pháp nào sau đây giúp phòng tránh cháy nổ trong phòng thực hành?
\choice
{Để dây điện chồng chéo và gần nguồn nhiệt}
{Cắm nhiều thiết bị công suất lớn vào cùng một ổ cắm}
{\True Kiểm tra thiết bị, dây dẫn và thu dọn vật dễ cháy trước khi thí nghiệm}
{Bật sẵn tất cả thiết bị trước khi giáo viên hướng dẫn}
\loigiai{
Kiểm tra thiết bị và thu dọn vật dễ cháy giúp giảm nguy cơ chập điện, quá tải và bắt lửa trong quá trình thực hành.
}
\end{ex}

% ===== Câu 182 =====
% ID: L10C1B2-96-TN
% Mức: NB
\begin{ex}
% ID: L10C1B2-96-TN
% Mức: NB
Khi thấy một người bị điện giật và vẫn đang tiếp xúc với nguồn điện, việc cần làm đầu tiên là gì?
\choice
{Chạy đến kéo nạn nhân bằng tay trần}
{\True Ngắt nguồn điện nếu có thể thực hiện an toàn}
{Cho nạn nhân uống nước ngay}
{Xoa dầu vào tay nạn nhân}
\loigiai{
Việc đầu tiên là phải tách nạn nhân khỏi nguồn điện bằng cách ngắt điện nếu an toàn. Không được chạm tay trần vào nạn nhân khi chưa ngắt điện.
}
\end{ex}

% ===== Câu 183 =====
% ID: L10C1B2-97-TN
% Mức: NB
\begin{ex}
% ID: L10C1B2-97-TN
% Mức: NB
Hành động nào sau đây là sai khi sơ cứu người bị điện giật?
\choice
{Ngắt nguồn điện trước khi tiếp xúc với nạn nhân}
{Gọi người hỗ trợ và gọi cấp cứu}
{\True Chạm tay trần kéo nạn nhân khi nạn nhân còn tiếp xúc với điện}
{Dùng vật khô cách điện để gạt dây điện}
\loigiai{
Chạm tay trần vào nạn nhân khi chưa tách khỏi nguồn điện có thể làm người cứu bị điện giật.
}
\end{ex}

% ===== Câu 184 =====
% ID: L10C1B2-98-TN
% Mức: TH
\begin{ex}
% ID: L10C1B2-98-TN
% Mức: TH
Bình khí carbonic ($\mathrm{CO}_2$) tuyệt đối không được sử dụng để dập tắt đám cháy có sự tham gia của kim loại hoạt động mạnh nào sau đây?
\choice
{Đồng ($\mathrm{Cu}$)}
{Sắt ($\mathrm{Fe}$)}
{\True Magnesium ($\mathrm{Mg}$)}
{Vàng ($\mathrm{Au}$)}
\loigiai{
Magnesium là kim loại hoạt động mạnh, ở nhiệt độ cao có khả năng khử khí $\mathrm{CO}_2$ mạnh mẽ tạo ra carbon và tỏa nhiệt lượng cực lớn: $2\mathrm{Mg} + \mathrm{CO}_2 \rightarrow 2\mathrm{MgO} + \mathrm{C}$. Phản ứng này tỏa nhiệt làm đám cháy bùng phát dữ dội hơn và muội carbon nóng đỏ sinh ra có thể gây nổ.
}
\end{ex}

% ===== Câu 185 =====
% ID: L10C1B2-99-TN
% Mức: TH
\begin{ex}
% ID: L10C1B2-99-TN
% Mức: TH
Khi trong phòng thực hành xảy ra đám cháy nhỏ từ cồn hoặc thiết bị điện nhưng không có sẵn bình chữa cháy chuyên dụng, vật liệu dễ kiếm nào sau đây dập lửa hiệu quả và an toàn nhất?
\choice
{Một tấm vải khô thưa}
{\True Một lượng cát khô đổ phủ kín lên ngọn lửa}
{Sách vở hoặc giấy báo khô gom lại ném vào đám cháy}
{Quạt giấy thổi mạnh vào ngọn lửa để khí độc phân tán}
\loigiai{
Cát khô là chất chữa cháy rất hiệu quả cho đám cháy điện hoặc cồn nhờ khả năng bao phủ bề mặt tốt, ngăn chặn triệt để sự tiếp xúc giữa chất cháy và oxygen trong không khí.
}
\end{ex}



% ID: Xuất từ robot vPT-Tác giả câu hỏi: CN
% Mức: TH
%=== Câu 1 ===%
\begin{ex}
% ID: Xuất từ robot vPT-Tác giả câu hỏi: CN
% Mức: TH %[Xuất từ robot vPT-Tác giả câu hỏi: CN]%[Mã câu: N10:45]
Đâu là cách đúng để đảm bảo an toàn khi sử dụng thiết bị điện trong phòng thí nghiệm Vật lí?
\choice
{Để dây điện tự do để dễ dàng kiểm tra}
{Sử dụng các thiết bị điện khi chưa được giáo viên cho phép}
{\True Chỉ cắm phích vào ổ cắm khi hiệu điện thế phù hợp}
{Không quan tâm đến kí hiệu và thông số trên thiết bị}
\loigiai{Để đảm bảo an toàn khi sử dụng thiết bị điện trong phòng thí nghiệm, cần tuân thủ các nguyên tắc sau:
- **Chỉ cắm phích vào ổ cắm khi hiệu điện thế phù hợp với thiết bị** để tránh quá tải, hỏng hóc hoặc cháy nổ. Đây là biện pháp an toàn cơ bản và quan trọng nhất.
- Luôn đọc và hiểu rõ các kí hiệu, thông số kĩ thuật trên thiết bị trước khi sử dụng để vận hành đúng cách.
- Chỉ sử dụng thiết bị điện khi có sự hướng dẫn hoặc cho phép của giáo viên để đảm bảo an toàn và tuân thủ quy định phòng thí nghiệm.
- Dây điện cần được sắp xếp gọn gàng, cố định, tránh vướng víu hoặc bị hư hỏng, tránh nguy cơ chập điện hoặc vấp ngã.}
\end{ex}

% ID: Xuất từ robot vPT-Tác giả câu hỏi: CN
% Mức: TH
%=== Câu 2 ===%
\begin{ex}
% ID: Xuất từ robot vPT-Tác giả câu hỏi: CN
% Mức: TH %[Xuất từ robot vPT-Tác giả câu hỏi: CN]%[Mã câu: N10:46]
Khi sử dụng các thiết bị điện trong phòng thí nghiệm Vật lí, việc cắm phích/giắc cắm vào ổ cắm nào sau đây là sai quy tắc an toàn?
\choice
{Ổ cắm mà dây điện không gây rối}
{Ổ cắm có hiệu điện thế phù hợp với hiệu điện thế định mức của dụng cụ}
{\True Ổ cắm gần nơi có thiết bị điện có nhiệt độ cao}
{Ổ cắm cách xa nơi có nước và các dung dịch dễ cháy}
\loigiai{Việc cắm phích/giắc cắm vào ổ cắm gần nơi có thiết bị điện có nhiệt độ cao là sai quy tắc an toàn. Nhiệt độ cao có thể làm hỏng lớp cách điện của dây điện, dẫn đến nguy cơ chập điện, đoản mạch hoặc cháy nổ. Các phương án còn lại đều là những quy tắc an toàn đúng đắn cần tuân thủ khi sử dụng thiết bị điện trong phòng thí nghiệm.}
\end{ex}

% ID: Xuất từ robot vPT-Tác giả câu hỏi: CN
% Mức: TH
%=== Câu 3 ===%
\begin{ex}
% ID: Xuất từ robot vPT-Tác giả câu hỏi: CN
% Mức: TH %[Xuất từ robot vPT-Tác giả câu hỏi: CN]%[Mã câu: N10:47]
Khi sử dụng thiết bị đo điện trong phòng thí nghiệm Vật lí, người sử dụng nên làm gì để đảm bảo độ chính xác của kết quả đo?
\choice
{Sử dụng dụng cụ bảo hộ chuyên dùng cho thiết bị này}
{Đo nhiệt độ của thiết bị trước khi tiến hành đo điện}
{\True Kiểm tra và hiệu chuẩn thiết bị định kỳ theo hướng dẫn của nhà sản xuất}
{Sử dụng nước để làm mát thiết bị trước khi đo}
\end{ex}

% ID: Xuất từ robot vPT-Tác giả câu hỏi: CN
% Mức: TH
%=== Câu 4 ===%
\begin{ex}
% ID: Xuất từ robot vPT-Tác giả câu hỏi: CN
% Mức: TH %[Xuất từ robot vPT-Tác giả câu hỏi: CN]%[Mã câu: N10:48]
Trong bài thực hành có sử dụng mạch điện nhưng khi lắp ráp xong mạch điện, báo cáo giáo viên phụ trách rồi cắm vào nguồn điện nhưng mạch không vào điện thì học sinh cần
\choice
{kiểm tra lại mạch điện}
{kiểm tra nguồn điện}
{\True ngắt mạch điện ra khỏi nguồn sau đó kiểm tra mạch điện và nguồn điện}
{ngắt mạch điện ra khỏi nguồn}
\end{ex}

% ID: Xuất từ robot vPT-Tác giả câu hỏi: CN
% Mức: TH
%=== Câu 5 ===%
\begin{ex}
% ID: Xuất từ robot vPT-Tác giả câu hỏi: CN
% Mức: TH %[Xuất từ robot vPT-Tác giả câu hỏi: CN]%[Mã câu: N10:49]
Người sử dụng phòng thí nghiệm Vật lí cần chú ý gì khi làm việc với các thiết bị có nhiệt độ cao?
\choice
{Sử dụng chất chống cháy cho môi trường làm việc}
{\True Đeo trang bị bảo hộ phù hợp}
{Sử dụng nước để làm mát các thiết bị}
{Đảm bảo có đủ bình chữa cháy gần đó}
\loigiai{Khi làm việc với các thiết bị có nhiệt độ cao, nguy cơ bỏng là rất lớn. Do đó, việc đeo trang bị bảo hộ cá nhân phù hợp như găng tay chịu nhiệt, kính bảo hộ và áo choàng thí nghiệm là biện pháp quan trọng hàng đầu để bảo vệ người sử dụng. Các phương án khác không phải là ưu tiên hàng đầu hoặc có thể không phù hợp (ví dụ, sử dụng nước để làm mát có thể gây nguy hiểm tùy thuộc vào loại thiết bị).}
\end{ex}

% ID: Xuất từ robot vPT-Tác giả câu hỏi: CN
% Mức: TH
%=== Câu 6 ===%
\begin{ex}
% ID: Xuất từ robot vPT-Tác giả câu hỏi: CN
% Mức: TH %[Xuất từ robot vPT-Tác giả câu hỏi: CN]%[Mã câu: N10:50]
Vì sao cần đeo kính bảo hộ khi sử dụng các thiết bị đun nóng và làm việc với các chất nóng?
\choice
{Để tránh bị đau khi sử dụng thiết bị}
{\True Để bảo vệ mắt khỏi các tia nhiệt và mắt bị nóng}
{Để tránh gây ra bệnh tật cho mắt}
{Để bảo vệ mắt khỏi bụi và dầu}
\loigiai{Khi sử dụng các thiết bị đun nóng và làm việc với các chất nóng, mắt có thể bị tổn thương do tiếp xúc trực tiếp với nhiệt độ cao (bỏng), các tia hồng ngoại phát ra từ nguồn nhiệt, hoặc do bắn tóe các chất nóng (như nước sôi, dầu nóng, hóa chất nóng chảy). Kính bảo hộ giúp tạo một hàng rào vật lý, bảo vệ mắt khỏi những nguy cơ này, đặc biệt là các tia nhiệt và nguy cơ bị bỏng hoặc tổn thương do chất nóng bắn vào.}
\end{ex}

% ID: Xuất từ robot vPT-Tác giả câu hỏi: CN
% Mức: TH
%=== Câu 7 ===%
\begin{ex}
% ID: Xuất từ robot vPT-Tác giả câu hỏi: CN
% Mức: TH %[Xuất từ robot vPT-Tác giả câu hỏi: CN]%[Mã câu: N10:51]
Theo quy tắc an toàn, khi sử dụng thiết bị nhiệt và thủy tinh trong phòng thí nghiệm Vật lí, bạn nên làm gì để tránh nguy cơ bị thương?
\choice
{\True Đeo kính bảo hộ khi làm việc}
{Dùng bàn tay trần để cầm thiết bị}
{Kiểm tra lại dụng cụ trước khi sử dụng}
{Sử dụng các chất liệu thay thế thủy tinh}
\loigiai{Khi làm việc với thiết bị nhiệt và thủy tinh trong phòng thí nghiệm, nguy cơ lớn nhất là các mảnh vỡ thủy tinh bắn vào mắt hoặc hóa chất bắn tóe gây tổn thương mắt. Do đó, việc đeo kính bảo hộ là biện pháp an toàn quan trọng hàng đầu và bắt buộc để bảo vệ đôi mắt. Kiểm tra dụng cụ trước khi sử dụng là một bước cần thiết nhưng không phải là biện pháp trực tiếp nhất để ngăn ngừa chấn thương khi thủy tinh vỡ đột ngột hoặc có sự cố nhiệt. Sử dụng vật liệu thay thế không phải là hành động cá nhân trong quá trình thao tác. Dùng tay trần để cầm thiết bị nhiệt hoặc thủy tinh sắc nhọn là hành vi nguy hiểm, dễ gây bỏng hoặc đứt tay.}
\end{ex}

% ID: Xuất từ robot vPT-Tác giả câu hỏi: CN
% Mức: TH
%=== Câu 8 ===%
\begin{ex}
% ID: Xuất từ robot vPT-Tác giả câu hỏi: CN
% Mức: TH %[Xuất từ robot vPT-Tác giả câu hỏi: CN]%[Mã câu: N10:52]
Tại sao không nên lau thiết bị quang học bằng vải thô hoặc giấy?
\choice
{Vải thô và giấy có thể gây nứt vỡ thiết bị quang học}
{Vải thô và giấy không phù hợp với mục đích vệ sinh thiết bị quang học}
{Vải thô và giấy thường bám bụi, gây làm mờ và sai lệch kết quả thí nghiệm}
{\True Vải thô và giấy dễ làm trầy xước mặt kính}
\loigiai{Vải thô và giấy, đặc biệt là loại dùng thông thường, thường có cấu trúc sợi cứng và thô ráp. Khi dùng để lau các bề mặt quang học nhạy cảm như ống kính, gương, hoặc màn hình, các sợi này có thể tạo ra các vết xước nhỏ li ti (micro-scratches) trên bề mặt, làm giảm chất lượng hình ảnh hoặc hiệu suất của thiết bị theo thời gian. Thay vào đó, nên sử dụng vải microfiber chuyên dụng hoặc khăn lau ống kính chuyên dụng mềm mại và sạch sẽ.}
\end{ex}

% ID: Xuất từ robot vPT-Tác giả câu hỏi: CN
% Mức: TH
%=== Câu 9 ===%
\begin{ex}
% ID: Xuất từ robot vPT-Tác giả câu hỏi: CN
% Mức: TH %[Xuất từ robot vPT-Tác giả câu hỏi: CN]%[Mã câu: N10:53]
Tại sao việc vệ sinh và bảo quản thiết bị quang học rất quan trọng trong phòng thí nghiệm Vật lí?
\choice
{\True Để bảo vệ thiết bị và duy trì hiệu suất thí nghiệm}
{Để tránh bị phạt vì không tuân thủ quy định về vệ sinh}
{Chỉ để tăng tuổi thọ của thiết bị}
{Để duy trì sự an toàn của người sử dụng}
\loigiai{Việc vệ sinh và bảo quản thiết bị quang học là cực kỳ quan trọng để đảm bảo chúng hoạt động chính xác, cho ra kết quả thí nghiệm đáng tin cậy. Bụi bẩn, dầu mỡ hoặc các vết bẩn khác trên thấu kính, gương có thể làm giảm độ truyền sáng, gây nhiễu loạn hoặc biến dạng hình ảnh, dẫn đến sai lệch kết quả. Đồng thời, việc bảo quản tốt cũng giúp kéo dài tuổi thọ của thiết bị, tiết kiệm chi phí thay thế và sửa chữa.}
\end{ex}

% ID: Xuất từ robot vPT-Tác giả câu hỏi: CN
% Mức: TH
%=== Câu 10 ===%
\begin{ex}
% ID: Xuất từ robot vPT-Tác giả câu hỏi: CN
% Mức: TH %[Xuất từ robot vPT-Tác giả câu hỏi: CN]%[Mã câu: N10:54]
Khi sử dụng các thiết bị quang học trong phòng thí nghiệm, người sử dụng cần chú ý đến điều gì để giảm thiểu sai lệch kết quả thí nghiệm?
\choice
{Đảm bảo thiết bị không bị mốc, xước}
{Sử dụng nước để làm sạch thiết bị}
{Đeo kính bảo hộ khi sử dụng}
{\True Thực hiện kiểm định hàng ngày cho thiết bị}
\loigiai{Lời giải:
Kiểm định (hay hiệu chuẩn) thiết bị quang học hàng ngày hoặc định kỳ là một biện pháp quan trọng để đảm bảo độ chính xác và tin cậy của các phép đo. Quá trình này giúp phát hiện và điều chỉnh các sai lệch do hao mòn, thay đổi môi trường hoặc các yếu tố khác, từ đó giảm thiểu đáng kể sai lệch trong kết quả thí nghiệm. Các phương án khác không trực tiếp tập trung vào việc giảm thiểu sai lệch kết quả thí nghiệm.
Phân tích các phương án:
*   Thực hiện kiểm định hàng ngày cho thiết bị: Đây là biện pháp trực tiếp và hiệu quả nhất để đảm bảo thiết bị hoạt động chính xác và giảm thiểu sai lệch kết quả theo thời gian.
*   Đảm bảo thiết bị không bị mốc, xước: Đây là điều kiện tiên quyết cho hoạt động của thiết bị quang học, nhưng nó liên quan đến tình trạng vật lý ban đầu hoặc bảo quản, không phải là một hành động kiểm soát sai lệch định kỳ như kiểm định. Nếu thiết bị đã bị mốc, xước thì dù có kiểm định cũng khó có kết quả chính xác.
*   Đeo kính bảo hộ khi sử dụng: Đây là quy tắc an toàn cá nhân, không liên quan đến việc giảm thiểu sai lệch kết quả thí nghiệm.
*   Sử dụng nước để làm sạch thiết bị: Việc làm sạch thiết bị quang học cần tuân thủ quy trình và sử dụng dung dịch chuyên dụng. Sử dụng nước thông thường có thể gây hỏng hóc hoặc để lại vết bẩn, làm tăng sai lệch, chứ không phải giảm thiểu.}
\end{ex}

% ID: Xuất từ robot vPT-Tác giả câu hỏi: CN
% Mức: TH
%=== Câu 11 ===%
\begin{ex}
% ID: Xuất từ robot vPT-Tác giả câu hỏi: CN
% Mức: TH %[Xuất từ robot vPT-Tác giả câu hỏi: CN]%[Mã câu: N10:55]
Khi phòng thí nghiệm thực hành xuất hiện cháy thì ta cần phải
\choice
{dùng nước dập đám cháy}
{ngắt nguồn điện, dùng nước dập đám cháy}
{\True ngắt điện, di chuyển các chất dễ cháy ra ngoài và chống cháy lan, cứu người, cứu tài sản, dập tắt đám cháy}
{chạy ra khỏi phòng, đi tìm thêm người đến dập đám cháy}
\loigiai{Khi xảy ra cháy trong phòng thí nghiệm, cần ưu tiên an toàn và ngăn chặn đám cháy lan rộng. Đầu tiên, ngắt nguồn điện để tránh nguy hiểm và ngăn cháy lan do điện. Sau đó, nhanh chóng di chuyển các vật liệu dễ cháy ra khỏi khu vực cháy. Luôn ưu tiên cứu người bị nạn, sau đó mới đến tài sản nếu an toàn. Cuối cùng, sử dụng các phương tiện chữa cháy phù hợp (bình chữa cháy, cát, chăn dập lửa, không phải nước cho mọi loại đám cháy) để dập tắt đám cháy. Các phương án khác không đầy đủ hoặc có thể gây nguy hiểm (ví dụ: dùng nước để dập đám cháy điện hoặc hóa chất).}
\end{ex}

% ID: Xuất từ robot vPT-Tác giả câu hỏi: CN
% Mức: TH
%=== Câu 12 ===%
\begin{ex}
% ID: Xuất từ robot vPT-Tác giả câu hỏi: CN
% Mức: TH %[Xuất từ robot vPT-Tác giả câu hỏi: CN]%[Mã câu: N10:56]
Trong trường hợp sử dụng các thiết bị điện có nguy cơ gây điện giật, người sử dụng nên tuân thủ các biện pháp an toàn nào sau đây để giảm thiểu nguy cơ này?
\choice
{Dùng nước để làm mát thiết bị điện}
{Đeo găng tay cao su khi sử dụng}
{\True Sử dụng thiết bị điện có cách ly tốt}
{Đảm bảo thiết bị điện luôn có 1 dây chống giật}
\loigiai{Để giảm thiểu nguy cơ điện giật khi sử dụng thiết bị điện, biện pháp quan trọng hàng đầu là đảm bảo thiết bị có cách ly tốt. Cách ly tốt (insulation) ngăn chặn dòng điện tiếp xúc trực tiếp với vỏ ngoài của thiết bị hoặc người sử dụng, từ đó giảm nguy cơ điện giật.
- Phương án "Đeo găng tay cao su khi sử dụng" là một biện pháp bảo hộ cá nhân, nhưng không phải là biện pháp an toàn cơ bản cho thiết bị.
- Phương án "Đảm bảo thiết bị điện luôn có 1 dây chống giật" (hay RCD/ELCB) là một hệ thống an toàn quan trọng của mạch điện nhằm ngắt nguồn khi có dòng rò, nhưng cách ly tốt là yếu tố an toàn cơ bản của bản thân thiết bị và là biện pháp phòng ngừa tiên quyết.
- Phương án "Dùng nước để làm mát thiết bị điện" là hoàn toàn sai và cực kỳ nguy hiểm, vì nước dẫn điện và sẽ làm tăng nguy cơ điện giật.
Do đó, sử dụng thiết bị có cách ly tốt là biện pháp an toàn hiệu quả nhất.}
\end{ex}

% ID: Xuất từ robot vPT-Tác giả câu hỏi: CN
% Mức: TH
%=== Câu 13 ===%
\begin{ex}
% ID: Xuất từ robot vPT-Tác giả câu hỏi: CN
% Mức: TH %[Xuất từ robot vPT-Tác giả câu hỏi: CN]%[Mã câu: N10:57]
Hoạt động nào sau đây không được làm sau khi kết thúc giờ thí nghiệm?
\choice
{Vệ sinh sạch sẽ phòng thí nghiệm}
{Sắp xếp gọn gàng các thiết bị và dụng cụ thí nghiệm}
{\True Để các thiết bị nối với nguồn điện giúp duy trì năng lượng}
{Bỏ chất thải thí nghiệm vào nới quy định}
\loigiai{Sau khi kết thúc giờ thí nghiệm, các thiết bị phải được ngắt kết nối khỏi nguồn điện để đảm bảo an toàn, tránh chập điện, cháy nổ hoặc tiêu thụ điện không cần thiết. Việc sắp xếp gọn gàng, xử lí chất thải và vệ sinh phòng thí nghiệm là những hoạt động bắt buộc phải làm để duy trì an toàn và vệ sinh phòng thí nghiệm.}
\end{ex}

% ID: Xuất từ robot vPT-Tác giả câu hỏi: CN
% Mức: TH
%=== Câu 14 ===%
\begin{ex}
% ID: Xuất từ robot vPT-Tác giả câu hỏi: CN
% Mức: TH %[Xuất từ robot vPT-Tác giả câu hỏi: CN]%[Mã câu: N10:58]
Khi tiến hành thí nghiệm với các thiết bị điện, tại sao cần chọn chính xác thang đo và chức năng trên đồng hồ đo điện?
\choice
{Để giảm thiểu nguy cơ cháy nổ trong phòng thí nghiệm}
{\True Để đảm bảo chính xác trong đo lường và tránh hỏng thiết bị đo}
{Vì đó là yêu cầu của nhà sản xuất thiết bị}
{Chỉ là để tuân thủ quy định, không ảnh hưởng nhiều đến kết quả thí nghiệm}
\loigiai{Chọn chính xác thang đo và chức năng trên đồng hồ đo điện là nguyên tắc cơ bản và quan trọng nhất khi sử dụng thiết bị đo điện. Việc này giúp đảm bảo rằng giá trị đo được là chính xác và đáng tin cậy. Đồng thời, nó cũng bảo vệ đồng hồ đo khỏi bị hư hỏng do quá tải hoặc sử dụng sai chức năng (ví dụ: đo điện áp ở chức năng đo dòng điện có thể gây ngắn mạch và hỏng thiết bị). Các lựa chọn khác không nêu bật được tầm quan trọng cốt lõi của việc này.}
\end{ex}

% ID: Xuất từ robot vPT-Tác giả câu hỏi: CN
% Mức: TH
%=== Câu 15 ===%
\begin{ex}
% ID: Xuất từ robot vPT-Tác giả câu hỏi: CN
% Mức: TH %[Xuất từ robot vPT-Tác giả câu hỏi: CN]%[Mã câu: N10:59]
Trong quá trình thí nghiệm, nếu bạn phát hiện ra thiết bị có dấu hiệu hỏng hoặc không hoạt động đúng cách, bạn nên làm gì trước khi tiếp tục thí nghiệm?
\choice
{Đưa thiết bị đến các trung tâm sửa chữa}
{Sử dụng các phương pháp sửa chữa tạm thời}
{Tiếp tục sử dụng bình thường và báo cáo sau khi thí nghiệm kết thúc}
{\True Chuyển sang sử dụng thiết bị thay thế ngay lập tức}
\loigiai{Khi phát hiện thiết bị có dấu hiệu hỏng hoặc không hoạt động đúng cách trong quá trình thí nghiệm, ưu tiên hàng đầu là đảm bảo an toàn và tính chính xác của kết quả.
*   Phương án "Tiếp tục sử dụng bình thường và báo cáo sau khi thí nghiệm kết thúc" hoặc "Sử dụng các phương pháp sửa chữa tạm thời" là rất nguy hiểm, có thể gây hỏng thiết bị thêm, gây tai nạn hoặc làm sai lệch kết quả thí nghiệm.
*   Phương án "Đưa thiết bị đến các trung tâm sửa chữa" là một bước xử lý thiết bị hỏng, nhưng không phải là điều bạn làm *ngay trước khi tiếp tục thí nghiệm* để hoàn thành công việc.
*   Phương án "Chuyển sang sử dụng thiết bị thay thế ngay lập tức" là hành động an toàn và hiệu quả nhất trong tình huống này để có thể tiếp tục thí nghiệm mà không làm ảnh hưởng đến an toàn và chất lượng dữ liệu. Điều này ngụ ý rằng bạn đã loại bỏ thiết bị hỏng khỏi quá trình sử dụng.}
\end{ex}

% ID: Xuất từ robot vPT-Tác giả câu hỏi: CN
% Mức: TH
%=== Câu 16 ===%
\begin{ex}
% ID: Xuất từ robot vPT-Tác giả câu hỏi: CN
% Mức: TH %[Xuất từ robot vPT-Tác giả câu hỏi: CN]%[Mã câu: N10:60]
Nguy cơ cháy nổ trong phòng thực hành Vật lí có thể được giảm bớt bằng cách nào?
\choice
{\True Bố trí dây điện gọn gàng, không bị vướng khi qua lại}
{Sử dụng nước và các dung dịch dễ cháy cùng lúc}
{Để các hóa chất dễ bay hơi trong chai mở nắp}
{Đốt các chất thải thí nghiệm bên trong phòng thực hành}
\loigiai{Lời giải:
Để giảm nguy cơ cháy nổ trong phòng thực hành Vật lí, cần tuân thủ các quy tắc an toàn.
- Phương án "Bố trí dây điện gọn gàng, không bị vướng khi qua lại" giúp tránh chập điện, hở mạch hoặc làm đổ thiết bị, vật liệu dễ cháy, từ đó giảm nguy cơ cháy nổ. Đây là phương án đúng.
- Phương án "Sử dụng nước và các dung dịch dễ cháy cùng lúc" là hành động nguy hiểm, có thể gây phản ứng hóa học hoặc làm lan rộng đám cháy nếu dung dịch dễ cháy bị đổ hoặc bắt lửa.
- Phương án "Để các hóa chất dễ bay hơi trong chai mở nắp" là rất nguy hiểm vì hơi hóa chất dễ bay hơi có thể tích tụ và bắt lửa, gây cháy nổ.
- Phương án "Đốt các chất thải thí nghiệm bên trong phòng thực hành" là hành động cấm kị vì có thể gây cháy lan, phát sinh khói độc và các chất nguy hiểm khác.}
\end{ex}

% ID: Xuất từ robot vPT-Tác giả câu hỏi: CN
% Mức: TH
%=== Câu 17 ===%
\begin{ex}
% ID: Xuất từ robot vPT-Tác giả câu hỏi: CN
% Mức: TH %[Xuất từ robot vPT-Tác giả câu hỏi: CN]%[Mã câu: N10:61]
Nguy cơ cháy nổ trong phòng thí nghiệm Vật lí thường xảy ra khi nào?
\choice
{Khi không có giáo viên hướng dẫn}
{\True Khi sử dụng thiết bị quá tải}
{Khi không có người đeo kính bảo hộ}
{Khi sử dụng các hóa chất dễ cháy}
\loigiai{Nguy cơ cháy nổ trong phòng thí nghiệm Vật lí chủ yếu liên quan đến việc sử dụng các thiết bị điện tử, điện từ, và các nguồn nhiệt. Khi các thiết bị này hoạt động quá công suất, hoặc khi hệ thống điện bị quá tải, có thể dẫn đến quá nhiệt, chập điện, và gây cháy nổ. Mặc dù một số thí nghiệm vật lí có thể liên quan đến chất dễ cháy (ví dụ: cồn, khí đốt), nhưng việc sử dụng thiết bị điện quá tải là nguyên nhân phổ biến và trực tiếp hơn cho cháy nổ trong môi trường phòng thí nghiệm vật lí.}
\end{ex}

% ID: Xuất từ robot vPT-Tác giả câu hỏi: CN
% Mức: TH
%=== Câu 18 ===%
\begin{ex}
% ID: Xuất từ robot vPT-Tác giả câu hỏi: CN
% Mức: TH %[Xuất từ robot vPT-Tác giả câu hỏi: CN]%[Mã câu: N10:62]
Tại sao cần lưu ý đặc biệt đối với sự an toàn khi sử dụng các vật liệu và thiết bị trong thí nghiệm Vật lí có nguy cơ gây phát nổ?
\choice
{\True Để giảm thiểu nguy cơ cháy nổ trong phòng thí nghiệm}
{Để đảm bảo tuân thủ quy định về an toàn}
{Vì các chất này dễ bị oxi hóa}
{Vì đó là yêu cầu của giáo viên}
\loigiai{Lời giải: Khi sử dụng các vật liệu và thiết bị có nguy cơ gây phát nổ trong thí nghiệm Vật lí, việc lưu ý đặc biệt về an toàn là cực kỳ cần thiết. Mục tiêu hàng đầu là giảm thiểu tối đa nguy cơ cháy nổ, bảo vệ tính mạng, sức khỏe của người làm thí nghiệm và những người có mặt trong phòng thí nghiệm, cũng như tránh hư hại tài sản.}
\end{ex}

% ID: Xuất từ robot vPT-Tác giả câu hỏi: CN
% Mức: TH
%=== Câu 19 ===%
\begin{ex}
% ID: Xuất từ robot vPT-Tác giả câu hỏi: CN
% Mức: TH %[Xuất từ robot vPT-Tác giả câu hỏi: CN]%[Mã câu: N10:63]
Chọn đáp án sai khi nói về những quy tắc an toàn trong phòng thí nghiệm Vật lí.
\choice
{\True Tắt công tắc nguồn thiết bị điện sau khi cắm hoặc tháo thiết bị điện}
{Kiểm tra cẩn thận thiết bị, phương tiện, dụng cụ thí nghiệm trước khi sử dụng}
{Chỉ tiến hành thí nghiệm khi được sự cho phép của giáo viên hướng dẫn thí nghiệm}
{Đọc kĩ hướng dẫn sử dụng thiết bị và quan sát các chỉ dẫn, các kí hiệu trên các thiết bị thí nghiệm}
\loigiai{**Lời giải:**
Trong phòng thí nghiệm, để đảm bảo an toàn, cần phải tắt công tắc nguồn thiết bị điện *trước khi* cắm hoặc tháo thiết bị điện. Việc tắt nguồn *sau khi* cắm hoặc tháo là không đúng quy tắc an toàn, có thể gây nguy hiểm điện giật hoặc hỏng hóc thiết bị.
Các phương án còn lại đều là các quy tắc an toàn đúng khi làm việc trong phòng thí nghiệm.}
\end{ex}

% ID: Xuất từ robot vPT-Tác giả câu hỏi: CN
% Mức: TH
%=== Câu 20 ===%
\begin{ex}
% ID: Xuất từ robot vPT-Tác giả câu hỏi: CN
% Mức: TH %[Xuất từ robot vPT-Tác giả câu hỏi: CN]%[Mã câu: N10:64]
Để đảm bảo an toàn khi sử dụng thiết bị điện trong phòng thí nghiệm Vật lí, bạn nên làm gì?
\choice
{Bật công tắc nguồn thiết bị điện ngay sau khi cắm phích}
{Cắm phích điện vào ổ cắm khi tay chưa khô}
{\True Quan sát kỹ kí hiệu và thông số trên thiết bị}
{Để dây điện chạy ngang qua lối đi để dễ quan sát}
\loigiai{Để đảm bảo an toàn khi sử dụng thiết bị điện trong phòng thí nghiệm Vật lí, điều quan trọng nhất là phải hiểu rõ về thiết bị. Việc quan sát kỹ kí hiệu và thông số (như điện áp, dòng điện, công suất) trên thiết bị giúp người dùng sử dụng đúng cách, tránh quá tải, hỏng hóc hoặc gây nguy hiểm.Các phương án sai:- Cắm phích điện vào ổ cắm khi tay chưa khô: Nước là chất dẫn điện, việc này có thể gây giật điện nguy hiểm.- Bật công tắc nguồn thiết bị điện ngay sau khi cắm phích: Nên kiểm tra lại các kết nối đã chắc chắn và đúng chưa trước khi bật nguồn để tránh sự cố.- Để dây điện chạy ngang qua lối đi để dễ quan sát: Việc này tạo ra chướng ngại vật, dễ gây vấp ngã và làm hỏng dây điện, dẫn đến nguy cơ chập điện hoặc hở mạch.}
\end{ex}

% ID: Xuất từ robot vPT-Tác giả câu hỏi: CN
% Mức: TH
%=== Câu 21 ===%
\begin{ex}
% ID: Xuất từ robot vPT-Tác giả câu hỏi: CN
% Mức: TH %[Xuất từ robot vPT-Tác giả câu hỏi: CN]%[Mã câu: N10:65]
Theo quy tắc an toàn, khi sử dụng thiết bị đo điện, bạn nên làm gì để tránh hỏng thiết bị?
\choice
{\True Lựa chọn chức năng và thang đo phù hợp}
{Chọn thang đo tùy ý}
{Rút phích điện khi dây điện hở}
{Chọn đồng hồ đo điện đa năng}
\loigiai{Để tránh hỏng thiết bị đo điện và đảm bảo an toàn, điều quan trọng nhất là phải lựa chọn đúng chức năng (đo điện áp, dòng điện, điện trở,...) và thang đo (giới hạn đo) phù hợp với giá trị cần đo. Việc chọn sai chức năng hoặc thang đo có thể dẫn đến quá tải và làm hỏng thiết bị.}
\end{ex}

% ID: Xuất từ robot vPT-Tác giả câu hỏi: CN
% Mức: TH
%=== Câu 22 ===%
\begin{ex}
% ID: Xuất từ robot vPT-Tác giả câu hỏi: CN
% Mức: TH %[Xuất từ robot vPT-Tác giả câu hỏi: CN]%[Mã câu: N10:66]
Trong phòng thí nghiệm Vật lí, việc cần phải kiểm tra cẩn thận thiết bị, dụng cụ trước khi sử dụng nhằm mục đích gì?
\choice
{Để chuẩn bị sẵn sàng cho các thí nghiệm tiếp theo}
{Để đảm bảo thiết bị không bị hỏng hóc trong quá trình sử dụng}
{\True Để đảm bảo an toàn cho người sử dụng và kết quả thí nghiệm}
{Để chắc chắn rằng thiết bị đang hoạt động tốt}
\loigiai{Việc kiểm tra cẩn thận thiết bị, dụng cụ trước khi sử dụng trong phòng thí nghiệm Vật lí nhằm hai mục đích chính: đảm bảo an toàn cho người sử dụng (tránh các sự cố như điện giật, cháy nổ, hỏng hóc gây nguy hiểm) và đảm bảo thiết bị hoạt động chính xác, từ đó cho ra kết quả thí nghiệm đáng tin cậy. Do đó, phương án "Để đảm bảo an toàn cho người sử dụng và kết quả thí nghiệm" là đầy đủ và chính xác nhất.}
\end{ex}

% ID: Xuất từ robot vPT-Tác giả câu hỏi: CN
% Mức: TH
%=== Câu 23 ===%
\begin{ex}
% ID: Xuất từ robot vPT-Tác giả câu hỏi: CN
% Mức: TH %[Xuất từ robot vPT-Tác giả câu hỏi: CN]%[Mã câu: N10:67]
Quy tắc an toàn nào trong phòng thí nghiệm Vật lí được xem là quan trọng nhất để đảm bảo an toàn cho người sử dụng và thiết bị?
\choice
{Đeo đồ bảo hộ khi làm việc}
{Đảm bảo các thiết bị đo đúng chuẩn}
{Báo cáo ngay lập tức với giáo viên khi có sự cố}
{\True Tuân thủ các chỉ dẫn và hướng dẫn của giáo viên}
\loigiai{Lời giải: Tuân thủ các chỉ dẫn và hướng dẫn của giáo viên là quy tắc an toàn quan trọng nhất vì nó bao gồm tất cả các biện pháp an toàn khác. Giáo viên là người có chuyên môn và kinh nghiệm, các chỉ dẫn của họ sẽ giúp học sinh thực hiện thí nghiệm một cách an toàn và hiệu quả, từ việc sử dụng thiết bị đúng cách, đeo đồ bảo hộ đến cách xử lý các tình huống khẩn cấp.}
\end{ex}

% ID: Xuất từ robot vPT-Tác giả câu hỏi: CN
% Mức: TH
%=== Câu 24 ===%
\begin{ex}
% ID: Xuất từ robot vPT-Tác giả câu hỏi: CN
% Mức: TH %[Xuất từ robot vPT-Tác giả câu hỏi: CN]%[Mã câu: N10:68]
Tại sao không nên để nước hoặc các dung dịch dễ cháy gần thiết bị điện trong phòng thí nghiệm Vật lí?
\choice
{Để tránh bị ẩm độ cao trong phòng thí nghiệm}
{\True Vì nước và dung dịch dễ cháy có thể gây ra cháy nổ}
{Vì nước có thể làm hỏng thiết bị điện}
{Để tránh bị điện giật khi làm việc với thiết bị}
\end{ex}

% ID: Xuất từ robot vPT-Tác giả câu hỏi: CN
% Mức: TH
%=== Câu 25 ===%
\begin{ex}
% ID: Xuất từ robot vPT-Tác giả câu hỏi: CN
% Mức: TH %[Xuất từ robot vPT-Tác giả câu hỏi: CN]%[Mã câu: N10:69]
Tại sao không nên tự ý vứt bỏ các chất thải từ các thí nghiệm trong phòng thí nghiệm Vật lí?
\choice
{Vì việc này phải tuân theo quy định về phòng cháy chữa cháy}
{Để tránh bị phạt tiền}
{Vì đó là việc của người quản lý phòng thí nghiệm}
{\True Để giảm thiểu rủi ro đối với môi trường và sức khỏe con người}
\loigiai{Các chất thải từ thí nghiệm trong phòng thí nghiệm Vật lí, đặc biệt là các hóa chất, vật liệu điện tử, pin, có thể chứa các thành phần độc hại hoặc nguy hiểm. Việc tự ý vứt bỏ có thể gây ô nhiễm môi trường (đất, nước, không khí) và ảnh hưởng trực tiếp đến sức khỏe con người (gây ngộ độc, bỏng, bệnh tật, v.v.). Do đó, cần tuân thủ các quy định về xử lí chất thải nguy hại để đảm bảo an toàn.}
\end{ex}

% ID: Xuất từ robot vPT-Tác giả câu hỏi: CN
% Mức: TH
%=== Câu 26 ===%
\begin{ex}
% ID: Xuất từ robot vPT-Tác giả câu hỏi: CN
% Mức: TH %[Xuất từ robot vPT-Tác giả câu hỏi: CN]%[Mã câu: N10:70]
Trong các hoạt động dưới đây, hoạt động nào tuân thủ theo nguyên tắc an toàn khi làm việc với các nguồn phóng xạ?
\choice
{Ăn uống, trang điểm trong phòng nơi chó chất phóng xạ}
{Tiếp xúc trực tiếp với chất phóng xạ}
{\True Sử dụng phương tiện phòng hộ cá nhân như quần áo phòng hộ, găng tay, mũ, áo chì}
{Đổ rác thải phóng xạ ra khu vực rác thải sinh hoạt}
\loigiai{Lời giải: Để đảm bảo an toàn khi làm việc với các nguồn phóng xạ, cần tuân thủ nghiêm ngặt các nguyên tắc an toàn. Việc sử dụng phương tiện phòng hộ cá nhân như quần áo phòng hộ, găng tay, mũ, áo chì là một trong những biện pháp quan trọng nhất để giảm thiểu sự tiếp xúc với bức xạ. Các hoạt động như ăn uống, trang điểm trong khu vực có chất phóng xạ, đổ rác thải phóng xạ ra khu vực rác thải sinh hoạt, hay tiếp xúc trực tiếp với chất phóng xạ đều là những hành vi vi phạm nguyên tắc an toàn và có thể gây nguy hiểm nghiêm trọng đến sức khỏe.}
\end{ex}

% ID: Xuất từ robot vPT-Tác giả câu hỏi: CN
% Mức: TH
%=== Câu 27 ===%
\begin{ex}
% ID: Xuất từ robot vPT-Tác giả câu hỏi: CN
% Mức: TH %[Xuất từ robot vPT-Tác giả câu hỏi: CN]%[Mã câu: N10:71]
Khi sử dụng các thiết bị trong phòng thí nghiệm Vật lí, kí hiệu “+” hoặc màu đỏ mang ý nghĩa là
\choice
{đầu ra}
{cực âm}
{đầu vào}
{\True cực dương}
\loigiai{Trong các thiết bị điện, điện tử, kí hiệu dấu “+” hoặc màu đỏ thường dùng để chỉ cực dương (nơi có điện thế cao hơn), còn dấu “-” hoặc màu đen/xanh thường dùng để chỉ cực âm (nơi có điện thế thấp hơn).}
\end{ex}

% ID: Xuất từ robot vPT-Tác giả câu hỏi: CN
% Mức: TH
%=== Câu 28 ===%
\begin{ex}
% ID: Xuất từ robot vPT-Tác giả câu hỏi: CN
% Mức: TH %[Xuất từ robot vPT-Tác giả câu hỏi: CN]%[Mã câu: N10:72]
Khi sử dụng các thiết bị trong phòng thí nghiệm Vật lí, kí hiệu ểượảểậếĩ \begin{center} \includegraphics[width=\linewidth, keepaspectratio]{DataPicture/438_0_28_De_12_9_2026_h12_41_34_1.png} \end{center}
\choice
{Tránh ánh nắng chiếu trực tiếp}
{Nhiệt độ cao}
{Nơi cấm lửa}
{\True Chất dễ cháy}
\loigiai{Kí hiệu hình ngọn lửa trong hình thoi đỏ là biểu tượng chuẩn quốc tế (Hệ thống hài hòa toàn cầu - GHS) dùng để chỉ các chất hoặc vật liệu dễ cháy (flammable substances).}
\end{ex}

% ID: Xuất từ robot vPT-Tác giả câu hỏi: CN
% Mức: TH
%=== Câu 29 ===%
\begin{ex}
% ID: Xuất từ robot vPT-Tác giả câu hỏi: CN
% Mức: TH %[Xuất từ robot vPT-Tác giả câu hỏi: CN]%[Mã câu: N10:73]
Khi sử dụng các thiết bị trong phòng thí nghiệm Vật lí, kí hiệu  \begin{center} \includegraphics[width=\linewidth, keepaspectratio]{DataPicture/438_0_29_De_12_9_2026_h12_41_34_1.png} \end{center}
\choice
{Lối thoát hiểm}
{Tránh gió trực tiếp}
{\True Nơi có chất phóng xạ}
{Nơi cấm sử dụng quạt}
\loigiai{Kí hiệu hình tam giác màu vàng với hình ba cánh quạt màu đen ở giữa là biểu tượng quốc tế cảnh báo về sự hiện diện của chất phóng xạ hoặc nguy hiểm phóng xạ.}
\end{ex}

% ID: Xuất từ robot vPT-Tác giả câu hỏi: CN
% Mức: TH
%=== Câu 30 ===%
\begin{ex}
% ID: Xuất từ robot vPT-Tác giả câu hỏi: CN
% Mức: TH %[Xuất từ robot vPT-Tác giả câu hỏi: CN]%[Mã câu: N10:74]
Khi sử dụng các thiết bị trong phòng thí nghiệm Vật lí, kí hiệu  \begin{center} \includegraphics[width=\linewidth, keepaspectratio]{DataPicture/438_0_30_De_12_9_2026_h12_41_34_1.png} \end{center}
\choice
{Lưu ý cẩn thận}
{\True Nơi nguy hiểm về điện}
{Cảnh báo tia laser}
{Cẩn thận sét đánh}
\loigiai{Kí hiệu hình tam giác màu vàng có hình tia sét màu đen là kí hiệu cảnh báo nguy hiểm điện giật hoặc nơi có điện áp cao. Đây là một biểu tượng an toàn phổ biến được sử dụng để chỉ ra các mối nguy hiểm liên quan đến điện.}
\end{ex}

% ID: Xuất từ robot vPT-Tác giả câu hỏi: CN
% Mức: TH
%=== Câu 31 ===%
\begin{ex}
% ID: Xuất từ robot vPT-Tác giả câu hỏi: CN
% Mức: TH %[Xuất từ robot vPT-Tác giả câu hỏi: CN]%[Mã câu: N10:75]
Trong một buổi thực hành Vật lý ở Trường THPT Kim Liên, học sinh C và học sinh D được yêu cầu tiến hành thí nghiệm đo cường độ dòng điện qua một điện trở bằng ampe kế. Trước khi bắt đầu, giáo viên đã nhắc nhở các học sinh về việc kiểm tra các thiết bị điện trước khi sử dụng. Học sinh C phát hiện rằng ampe kế có dấu hiệu hỏng và kim chỉ thị không di chuyển về vị trí 0 khi không có dòng điện. Học sinh C báo cáo tình trạng này cho giáo viên và được hướng dẫn sử dụng một ampe kế khác. Trong khi đó, học sinh D kết nối mạch điện với nguồn điện, nhưng đã vô tình để dây dẫn tiếp xúc với nước trên bàn thí nghiệm. Học sinh D ngay lập tức tắt nguồn điện và lau khô bàn thí nghiệm trước khi tiếp tục tiến hành thí nghiệm.
\choiceTF
{\True  Học sinh C đã đúng khi báo cáo tình trạng hỏng của ampe kế cho giáo viên.
}
{ Việc để dây dẫn tiếp xúc với nước trên bàn thí nghiệm là an toàn nếu dây dẫn không bị hở.
}
{ Học sinh D đã thực hiện sai khi tắt nguồn điện và lau khô bàn thí nghiệm trước khi tiếp tục tiến hành thí nghiệm.
}
{ Học sinh không cần kiểm tra các thiết bị điện trước khi sử dụng nếu chúng đã được sử dụng trong các buổi thí nghiệm trước đó.
}
\loigiai{Trong phòng thí nghiệm, việc tuân thủ các quy tắc an toàn là vô cùng quan trọng để đảm bảo an toàn cho bản thân và những người xung quanh, đồng thời bảo vệ thiết bị.
- **Phương án 1 (Học sinh C):** Việc kiểm tra và báo cáo thiết bị hỏng là hành động đúng đắn và cần thiết. Ampe kế bị hỏng có thể gây ra kết quả sai hoặc thậm chí nguy hiểm nếu tiếp tục sử dụng. Đây là một quy tắc cơ bản trong việc sử dụng thiết bị thí nghiệm.
- **Phương án 2 (Dây dẫn tiếp xúc với nước):** Nước là chất dẫn điện. Ngay cả khi dây dẫn không bị hở, việc để dây dẫn điện tiếp xúc với nước có nguy cơ gây chập mạch, hỏa hoạn hoặc điện giật. Đây là hành vi không an toàn và vi phạm nghiêm trọng quy tắc an toàn điện trong phòng thí nghiệm.
- **Phương án 3 (Học sinh D xử lý sự cố):** Khi có sự cố như nước tiếp xúc với mạch điện, hành động đúng đắn nhất là ngắt nguồn điện ngay lập tức và làm khô khu vực bị ảnh hưởng trước khi tiếp tục. Điều này giúp ngăn ngừa nguy cơ điện giật và hỏng hóc thiết bị.
- **Phương án 4 (Kiểm tra thiết bị):** Luôn cần kiểm tra các thiết bị điện trước mỗi lần sử dụng, bất kể chúng đã được sử dụng trước đó hay chưa. Thiết bị có thể bị hỏng hóc hoặc có vấn đề không nhìn thấy được giữa các lần sử dụng, tiềm ẩn nguy hiểm. Đây là một nguyên tắc cơ bản để đảm bảo an toàn và độ chính xác của thí nghiệm.
Tóm lại, cả học sinh C và D đều đã có những hành động thể hiện sự hiểu biết và tuân thủ các nguyên tắc an toàn cơ bản trong phòng thí nghiệm.(Đúng)  Học sinh C đã đúng khi báo cáo tình trạng hỏng của ampe kế cho giáo viên.
(Vì):  Đúng(Sai)  Việc để dây dẫn tiếp xúc với nước trên bàn thí nghiệm là an toàn nếu dây dẫn không bị hở.
(Vì):  Sai vì không tuân theo quy tắc an toàn trong phòng Thí nghiệm(Sai)  Học sinh D đã thực hiện sai khi tắt nguồn điện và lau khô bàn thí nghiệm trước khi tiếp tục tiến hành thí nghiệm.
(Vì):  Sai(Sai)  Học sinh không cần kiểm tra các thiết bị điện trước khi sử dụng nếu chúng đã được sử dụng trong các buổi thí nghiệm trước đó.
(Vì):  Sai vì vi phạm quy tắc an toàn trong phòng Thí nghiệm}
\end{ex}

% ID: Xuất từ robot vPT-Tác giả câu hỏi: CN
% Mức: TH
%=== Câu 32 ===%
\begin{ex}
% ID: Xuất từ robot vPT-Tác giả câu hỏi: CN
% Mức: TH %[Xuất từ robot vPT-Tác giả câu hỏi: CN]%[Mã câu: N10:76]
Trong một buổi thực hành Vật lý tại Trường THPT Lý Thường Kiệt, học sinh G và học sinh H được yêu cầu tiến hành thí nghiệm đo lực từ của một nam châm. Giáo viên đã nhắc nhở các học sinh không để nam châm gần các thiết bị điện tử cá nhân như điện thoại di động hoặc đồng hồ. Học sinh G vô tình để nam châm gần điện thoại di động của mình và nhận thấy màn hình điện thoại bị nhiễu. Học sinh G lập tức di chuyển nam châm ra xa và kiểm tra lại các thiết bị khác để đảm bảo không có sự cố nào khác xảy ra. Trong khi đó, học sinh H, khi thao tác với nam châm, đã để nam châm rơi xuống và va chạm mạnh với bàn thí nghiệm. Học sinh H báo cáo sự cố này cho giáo viên và được hướng dẫn kiểm tra lại tình trạng của nam châm trước khi tiếp tục thí nghiệm.
\choiceTF
{ Học sinh G đã thực hiện sai khi di chuyển nam châm ra xa điện thoại di động khi phát hiện nhiễu.
}
{ Việc báo cáo sự cố nam châm rơi và va chạm mạnh với bàn thí nghiệm của học sinh H là không cần thiết.
}
{ Nam châm có thể để gần các thiết bị điện tử cá nhân mà không gây ra bất kỳ sự cố nào.
}
{ Học sinh không cần kiểm tra lại tình trạng của nam châm sau khi nó bị rơi và va chạm mạnh.
}
\loigiai{Câu hỏi này kiểm tra kiến thức và kỹ năng của học sinh về an toàn trong phòng thí nghiệm khi làm việc với nam châm. Việc tuân thủ hướng dẫn của giáo viên, xử lý đúng đắn các sự cố (như di chuyển nam châm ra xa thiết bị điện tử khi có nhiễu, hoặc báo cáo khi nam châm bị va chạm) là rất quan trọng để đảm bảo an toàn cho bản thân, thiết bị và duy trì hiệu quả thí nghiệm.(Sai)  Học sinh G đã thực hiện sai khi di chuyển nam châm ra xa điện thoại di động khi phát hiện nhiễu.
(Vì):   Sai, vì việc di chuyển nam châm ra xa là hành động đúng để tránh ảnh hưởng đến thiết bị điện tử và đảm bảo an toàn cho điện thoại.(Sai)  Việc báo cáo sự cố nam châm rơi và va chạm mạnh với bàn thí nghiệm của học sinh H là không cần thiết.
(Vì):   Sai, vì báo cáo sự cố là cần thiết để giáo viên nắm được tình hình, kiểm tra tình trạng của nam châm và đảm bảo an toàn trong phòng thí nghiệm.(Sai)  Nam châm có thể để gần các thiết bị điện tử cá nhân mà không gây ra bất kỳ sự cố nào.
(Vì):   Sai vì nam châm sẽ gây ảnh hưởng đến các thiết bị điện tử đặt cạnh nó(Sai)  Học sinh không cần kiểm tra lại tình trạng của nam châm sau khi nó bị rơi và va chạm mạnh.
(Vì):   Sai vi phạm quy tắc an toàn trong phòng Thí nghiệm, nam châm khi bị rơi và va chạm mạnh có thể bị ảnh hưởng đến từ tính của nó}
\end{ex}

% ID: Xuất từ robot vPT-Tác giả câu hỏi: CN
% Mức: TH
%=== Câu 33 ===%
\begin{ex}
% ID: Xuất từ robot vPT-Tác giả câu hỏi: CN
% Mức: TH %[Xuất từ robot vPT-Tác giả câu hỏi: CN]%[Mã câu: N10:77]
Trong một buổi thực hành Vật lý tại Trường THPT Hà Thành, học sinh I và học sinh J được yêu cầu tiến hành thí nghiệm xác định hệ số ma sát của một bề mặt bằng cách kéo một vật qua bề mặt đó. Giáo viên đã yêu cầu các học sinh sử dụng cân để đo khối lượng của vật trước khi tiến hành thí nghiệm. Học sinh I sử dụng cân để đo khối lượng của vật nhưng nhận thấy cân không ổn định và chỉ số liên tục thay đổi. Học sinh I báo cáo tình trạng này cho giáo viên và được hướng dẫn sử dụng một cân khác. Trong khi đó, học sinh J, khi kéo vật qua bề mặt, đã nhận thấy bề mặt có dầu mỡ làm giảm lực ma sát. Học sinh J tạm dừng thí nghiệm, làm sạch bề mặt và tiếp tục thí nghiệm sau khi đảm bảo bề mặt không còn dầu mỡ.
\choiceTF
{\True  Học sinh I đã đúng khi báo cáo tình trạng không ổn định của cân cho giáo viên.
}
{ Việc làm sạch bề mặt có dầu mỡ của học sinh J là không cần thiết vì lực ma sát không bị ảnh hưởng.
}
{ Học sinh không cần kiểm tra các thiết bị đo trước khi sử dụng trong thí nghiệm.
}
{ Khi phát hiện bề mặt có dầu mỡ, học sinh nên tiếp tục thí nghiệm mà không cần làm sạch.
}
\loigiai{Đây là một câu hỏi kiểm tra kỹ năng thực hành và tư duy khoa học trong thí nghiệm.
1.  Hành động của học sinh I: Khi cân không ổn định, việc báo cáo giáo viên và yêu cầu sử dụng thiết bị khác là hoàn toàn đúng đắn. Điều này đảm bảo tính chính xác của dữ liệu ban đầu (khối lượng vật), một yếu tố quan trọng trong việc xác định hệ số ma sát. Đây là một nguyên tắc cơ bản trong mọi thí nghiệm: đảm bảo thiết bị đo hoạt động chính xác.
2.  Hành động của học sinh J: Việc bề mặt có dầu mỡ sẽ ảnh hưởng đến lực ma sát (cụ thể là làm giảm lực ma sát). Để xác định hệ số ma sát của *bề mặt đó* trong điều kiện chuẩn (không có chất bôi trơn), việc làm sạch bề mặt là cần thiết. Tiếp tục thí nghiệm khi bề mặt bị ảnh hưởng bởi dầu mỡ sẽ cho kết quả sai lệch và không phản ánh đúng hệ số ma sát cần tìm. Hành động tạm dừng và làm sạch của học sinh J thể hiện sự cẩn trọng và hiểu biết về các yếu tố ảnh hưởng đến kết quả thí nghiệm.
Tóm lại, cả hai học sinh I và J đều có những hành động đúng đắn và cần thiết để đảm bảo tính chính xác và tin cậy của kết quả thí nghiệm. Các phương án đúng phản ánh những nguyên tắc cơ bản của thực hành khoa học: kiểm tra thiết bị, loại bỏ các yếu tố gây nhiễu, và báo cáo vấn đề khi cần thiết.(Đúng)  Học sinh I đã đúng khi báo cáo tình trạng không ổn định của cân cho giáo viên.
(Vì):   Đúng.(Sai)  Việc làm sạch bề mặt có dầu mỡ của học sinh J là không cần thiết vì lực ma sát không bị ảnh hưởng.
(Vì):   Sai vì lực ma sát sẽ giảm nếu bề mặt có dầu mỡ, làm sai lệch kết quả thí nghiệm.(Sai)  Học sinh không cần kiểm tra các thiết bị đo trước khi sử dụng trong thí nghiệm.
(Vì):   Đúng.(Sai)  Khi phát hiện bề mặt có dầu mỡ, học sinh nên tiếp tục thí nghiệm mà không cần làm sạch.
(Vì):   Sai vì gây ảnh hưởng đến kết quả thí nghiệm.}
\end{ex}

% ID: Xuất từ robot vPT-Tác giả câu hỏi: CN
% Mức: TH
%=== Câu 34 ===%
\begin{ex}
% ID: Xuất từ robot vPT-Tác giả câu hỏi: CN
% Mức: TH %[Xuất từ robot vPT-Tác giả câu hỏi: CN]%[Mã câu: N10:78]
Trong một buổi thực hành Vật lý, học sinh E và học sinh F được yêu cầu tiến hành thí nghiệm xác định điện trở của một dây dẫn bằng phương pháp cầu Wheatstone. Trước khi bắt đầu, giáo viên đã yêu cầu các học sinh đeo kính bảo hộ và găng tay bảo hộ. Học sinh E kết nối các dây dẫn và nhận thấy một trong các điện trở có dấu hiệu nóng lên quá mức. Học sinh E lập tức ngắt nguồn điện và báo cáo cho giáo viên. Học sinh F, trong khi kiểm tra các thiết bị, phát hiện rằng nguồn điện có dấu hiệu chập chờn. Học sinh F thay thế nguồn điện bằng một nguồn điện mới và kiểm tra lại toàn bộ mạch điện trước khi bắt đầu thí nghiệm.
\choiceTF
{ Học sinh E đã sai khi ngắt nguồn điện và báo cáo cho giáo viên khi phát hiện điện trở nóng lên quá mức, mà lẽ ra nên tự tìm cách khắc phục.
}
{\True  Việc thay thế nguồn điện chập chờn bằng một nguồn điện mới của học sinh F là cần thiết và đúng đắn để đảm bảo an toàn và kết quả thí nghiệm chính xác.
}
{ Học sinh không cần thiết đeo kính bảo hộ và găng tay bảo hộ khi làm việc với các thiết bị điện nếu chỉ là các mạch điện thế thấp.
}
{ Khi phát hiện thiết bị điện có dấu hiệu bất thường, học sinh có thể tự xử lý hoặc tiếp tục thí nghiệm mà không cần báo cáo ngay cho giáo viên nếu cảm thấy không quá nguy hiểm.
}
\loigiai{Câu hỏi này đánh giá kiến thức về an toàn trong phòng thí nghiệm.
1.  **Hành động của học sinh E:** Khi phát hiện điện trở nóng lên quá mức, việc ngắt nguồn điện ngay lập tức và báo cáo giáo viên là hành động hoàn toàn đúng đắn. Điều này giúp ngăn chặn nguy cơ cháy nổ, hỏng hóc thiết bị và đảm bảo an toàn cho người thực hiện.
2.  **Hành động của học sinh F:** Nguồn điện chập chờn sẽ dẫn đến kết quả đo không chính xác và tiềm ẩn nguy hiểm về điện. Do đó, việc thay thế nguồn điện mới là cần thiết và thể hiện sự cẩn trọng.
3.  **Quy định về kính bảo hộ và găng tay:** Trong phòng thí nghiệm, đặc biệt với các thí nghiệm liên quan đến điện, việc đeo kính bảo hộ và găng tay bảo hộ là quy tắc an toàn bắt buộc. Chúng bảo vệ mắt và tay khỏi các rủi ro như tia lửa điện, nóng chảy vật liệu, hoặc các sự cố khác.
4.  **Báo cáo giáo viên:** Mọi dấu hiệu bất thường của thiết bị hay mạch điện đều phải được báo cáo ngay lập tức cho giáo viên để nhận được hướng dẫn xử lý an toàn và kịp thời.(Sai)  Học sinh E đã sai khi ngắt nguồn điện và báo cáo cho giáo viên khi phát hiện điện trở nóng lên quá mức, mà lẽ ra nên tự tìm cách khắc phục.
(Vì):  Sai. Việc ngắt nguồn điện và báo cáo giáo viên là hành động đúng đắn và cần thiết để đảm bảo an toàn cho người thực hiện thí nghiệm, tránh hư hỏng thiết bị và nguy cơ cháy nổ. Tự ý khắc phục có thể gây nguy hiểm.(Đúng)  Việc thay thế nguồn điện chập chờn bằng một nguồn điện mới của học sinh F là cần thiết và đúng đắn để đảm bảo an toàn và kết quả thí nghiệm chính xác.
(Vì):  Đúng. Nguồn điện không ổn định sẽ làm ảnh hưởng nghiêm trọng đến kết quả thí nghiệm, ngoài ra còn gây nguy hiểm cho người sử dụng và thiết bị nếu cố tình sử dụng. Do đó, việc thay thế là cần thiết.(Sai)  Học sinh không cần thiết đeo kính bảo hộ và găng tay bảo hộ khi làm việc với các thiết bị điện nếu chỉ là các mạch điện thế thấp.
(Vì):  Sai. Ngay cả với mạch điện thế thấp, việc đeo kính bảo hộ và găng tay bảo hộ vẫn là quy tắc an toàn cơ bản trong phòng thí nghiệm để bảo vệ mắt và tay khỏi các sự cố không lường trước.(Sai)  Khi phát hiện thiết bị điện có dấu hiệu bất thường, học sinh có thể tự xử lý hoặc tiếp tục thí nghiệm mà không cần báo cáo ngay cho giáo viên nếu cảm thấy không quá nguy hiểm.
(Vì):  Sai. Mọi dấu hiệu bất thường của thiết bị hay mạch điện đều phải được báo cáo ngay lập tức cho giáo viên. Tự ý xử lý hoặc tiếp tục thí nghiệm có thể dẫn đến nguy hiểm và hỏng hóc thiết bị.}
\end{ex}

